\chapter{Setting (Sections 2 and 3)}
\label{chap:setting}

Throughout, $\X \subset \R^d$ is a finite set of arms ($\R^d$ is \texttt{EuclideanSpace} over a finite type $\iota$, with
$d = \abs{\iota}$) and $T \in \N$ is the budget. Rounds are numbered $0, \dots, T - 1$.

\begin{definition}[Non-stationary linear bandit]
  \label{def:setting}
  \lean{MaynardZhang2026Complexity.avgParam, Bandits.Linear.nonstationaryLinearEnv}
  \leanok
  An instance is a parameter sequence $\theta = (\theta_t)_{t \in \N}$ in $\R^d$ and noise laws
  $(\xi_t)_t$ on $\R$: at round $t$ the learner plays $x_t \in \X$ and observes
  $r_t = x_t^\top \theta_t + \eps_t$ with $\eps_t \sim \xi_t$ independent of the past
  (\texttt{nonstationaryLinearEnv}, an LML environment whose reward kernel at round $t$ is
  the linear kernel of $\theta_t$). The noise is centered $1$-sub-Gaussian in the upper bound and
  standard Gaussian in the lower bound. The average parameter is
  $\thetabar_T = \frac1T \sum_{t < T} \theta_t$.
\end{definition}

\begin{definition}[Extreme points and adjacency, Definition 1]
  \label{def:adjacency}
  \lean{Set.vertices, Set.IsAdjacent, Set.adjacentPairs, Set.adjacentTo, Set.vertexPairs}
  \leanok
  The \emph{vertices} $\V$ of $\X$ are the exposed points of $\conv \X$. Two distinct vertices
  $x, x'$ are \emph{adjacent} if the segment $[x, x']$ is an exposed subset of $\conv \X$ (an edge).
  $\Y = \{(x, x') \in \V^2 : x \ne x'\}$ is the set of pairs of distinct vertices,
  $\I \subseteq \Y$ the set of adjacent pairs, and $\I^x$ the set of vertices adjacent to $x$.
\end{definition}

\begin{definition}[Best arm and min-gap]
  \label{def:min_gap}
  \lean{MaynardZhang2026Complexity.IsBestArm, MaynardZhang2026Complexity.minGap,
    MaynardZhang2026Complexity.MinGapGE}
  \leanok
  \uses{def:setting, def:adjacency}
  An arm $x \in \X$ is a \emph{best arm} of $\theta$ over $T$ rounds if it maximizes
  $y \mapsto y^\top \thetabar_T$ over $\X$. The \emph{min-gap} of the best arm $x_\star$ is
  $\Delta_{(1)} = \min_{x \in \V \setminus \{x_\star\}} (x_\star - x)^\top \thetabar_T$. The sequence
  $\theta$ has min-gap at least $\Delta$ if some $x \in \X$ satisfies $(x - y)^\top \thetabar_T \ge
  \Delta$ for every vertex $y \ne x$; for $\Delta > 0$ this $x$ is then the unique best arm
  (Lemma~\ref{lem:min_gap_best_arm}).
\end{definition}

\begin{definition}[Feasible sequences, Definition 2]
  \label{def:feasible}
  \lean{MaynardZhang2026Complexity.FeasibleSeqs}
  \leanok
  \uses{def:min_gap}
  Two parameter sequences are \emph{feasible} for $(\X, \Delta)$ if they have different best arms
  and each has min-gap at least $\Delta$.
\end{definition}

\begin{definition}[Designs and Mahalanobis norms]
  \label{def:design}
  \lean{Learning.IsDesign, Learning.designMatrix, Matrix.mahalanobisSq,
    MaynardZhang2026Complexity.PosDefDesign}
  \leanok
  A \emph{design} on $\X$ is a probability vector $\lambda \in \triangle_{\X}$ (a finitely supported
  function with positive weights on $\X$ summing to $1$); its design matrix is
  $A(\lambda) = \sum_{x} \lambda_x x x^\top$. For a matrix $A$, $\norm{v}_A^2 = v^\top A v$.
  \texttt{PosDefDesign} is the set of designs with positive definite design matrix.
\end{definition}

\begin{definition}[Adjacency complexity]
  \label{def:h_adjacent}
  \lean{MaynardZhang2026Complexity.adjValue, MaynardZhang2026Complexity.adjDesignValue,
    MaynardZhang2026Complexity.hAdjacent}
  \leanok
  \uses{def:adjacency, def:design}
  \[
    \Hadj(\X, \Delta) = \frac{1}{\Delta^2} \min_{\lambda} \max_{(x, x') \in \I}
      \norm{x - x'}^2_{A(\lambda)^{-1}},
  \]
  the minimum being over the designs with positive definite design matrix (an infimum in Lean,
  \texttt{adjDesignValue}; the inner maximum is \texttt{adjValue}).
\end{definition}

\begin{lemma}[Adjacency lemma, Lemma 1]
  \label{lem:adjacency}
  \lean{MaynardZhang2026Complexity.exists_pos_inner_iff}
  \leanok
  \uses{def:adjacency, lem:local_global}
  Let $x \in \V$ and $\theta \in \R^d$. There is $y \in \X$ with $(y - x)^\top \theta > 0$ if and
  only if there is $z \in \I^x$ with $(z - x)^\top \theta > 0$.
\end{lemma}

\begin{proof}
  \leanok
  \uses{lem:vertices_subset, lem:local_global}
  Backward: adjacent vertices are points of $\X$ (Lemma~\ref{lem:vertices_subset}). Forward:
  Lemma~\ref{lem:local_global}.
\end{proof}

\begin{remark}
  The paper deduces Lemma~\ref{lem:adjacency} from the cone property of polytopes,
  $\conv \X \subseteq x + \cone\{z - x : z \in \I^x\}$ (\cite[Lemma 3.6]{ziegler1995lectures}),
  which is not in Mathlib. The blueprint proves the lemma directly
  (Lemma~\ref{lem:local_global}) and deduces the cone property in the form needed for
  Lemma~\ref{lem:ratio_adjacent} (Lemma~\ref{lem:cone_form}).
\end{remark}

\begin{lemma}[The arm of the min-gap is the best arm]
  \label{lem:min_gap_best_arm}
  \lean{Set.Finite.mem_vertices_of_forall_vertices}
  \leanok
  \uses{def:min_gap, lem:strict_expose, lem:max_face}
  Let $\Delta > 0$ and let $x \in \X$ beat every other vertex by $\Delta$ for $\thetabar_T$. Then $x$
  is a vertex, $x^\top \thetabar_T > y^\top \thetabar_T$ for all $y \in \X \setminus \{x\}$, and $x$
  is the unique best arm.
\end{lemma}

\begin{proof}
  \leanok
  \uses{lem:strict_expose, lem:max_face}
  Every vertex $v$ has $v^\top \thetabar_T \le x^\top \thetabar_T$, with equality only for $v = x$.
  Let $y \in \X$ with $y^\top \thetabar_T \ge x^\top \thetabar_T$. As $y \in \conv \X = \conv \V$
  (Lemma~\ref{lem:strict_expose}), $y$ is a convex combination of the vertices $v$ with
  $v^\top \thetabar_T = x^\top \thetabar_T$ (Lemma~\ref{lem:max_face}), of which there is at most
  one, $x$: so this set is $\{x\}$, $x$ is a vertex and $y = x$. Applied to $y = x$ this shows
  that $x$ is a vertex, and applied to $y \ne x$ that $y^\top \thetabar_T < x^\top \thetabar_T$.
\end{proof}
